% GENERATED FILE. Edit canonical lessons, then run npm run generate:books.
% canonical-source-hash: fnv1a64:2aa45d94722cfaba
% canonical-lessons: MW-C144-sundar, MW-C144-rupalo, MW-C144-javan, MW-C144-dubalo, MW-C144-thigno

\chapter{Beautiful, Good-looking, Young, Thin, Short}
\label{ch:mw-beautiful-good-looking-young-thin-short}
\hlchaptermodality{144}

\begin{chapteropening}
\textbf{By the end of this chapter:} \emph{I can say beautiful, good-looking, young, thin and short.}

Complete the last lesson of chapter 144: I can say beautiful, good-looking, young, thin and short.
\end{chapteropening}

\section[\texorpdfstring{\mw{सुंदर} (sundar)}{सुंदर (sundar)}]{\mw{सुंदर} (sundar) — beautiful}

\label{lesson:MW-C144-sundar}



\begin{quote}
Before the new one: say the Marwadi for bitter, then the Marwadi for spicy.
\end{quote}

\subsection*{You'll want to know: \mw{सुंदर}}
\textbf{\mw{सुंदर}} — \emph{sundar} — \textquotedblleft{}beautiful\textquotedblright{}.

\begin{cousinweb}
\textbf{\mw{सुंदर}} (\emph{sundar}) is Sanskrit \emph{sundara}, \textquotedblleft{}beautiful\textquotedblright{}, taken over whole as a learned loan, a tatsama (CDIAL 13474). CDIAL asks whether the Sanskrit word is itself a Middle Indo-Aryan form of \emph{sūnara}, \textquotedblleft{}excellent\textquotedblright{}.
\end{cousinweb}

\begin{grammarlens}[title={What you've built}]
One more word for this chapter.
\end{grammarlens}

\subsection*{Your turn}
\begin{itemize}
\raggedright
  \item \emph{Say it:} \emph{sundar}
  \item \emph{Say it:} \emph{sundar}, once more
  \item \emph{Say it:} say \emph{sundar}
  \item \emph{Recall:} say the Marwadi for bitter, then the Marwadi for spicy, then say \emph{sundar} again
\end{itemize}

\begin{tcolorbox}[breakable,colback=teal!4,colframe=teal!35!black,title={Before you move on}]
What does \mw{सुंदर} mean? (\textquotedblleft{}Beautiful\textquotedblright{}.) Say it once more.
\end{tcolorbox}
\section[\texorpdfstring{\mw{रूपाळो} (rūpāḷo)}{रूपाळो (rūpāḷo)}]{\mw{रूपाळो} (rūpāḷo) — good-looking}

\label{lesson:MW-C144-rupalo}



\begin{quote}
Before the new one: say the Marwadi for spicy, then the Marwadi for beautiful.
\end{quote}

\subsection*{You'll want to know: \mw{रूपाळो}}
\textbf{\mw{रूपाळो}} — \emph{rūpāḷo} — \textquotedblleft{}good-looking\textquotedblright{}.

\begin{cousinweb}
\textbf{\mw{रूपाळो}} (\emph{rūpāḷo}) is built with the suffix \emph{-āḷo} on \emph{rūp}, \textquotedblleft{}beauty\textquotedblright{}, from Sanskrit \emph{rūpa}, \textquotedblleft{}form, beauty\textquotedblright{}. CDIAL lists Hindi \emph{rūṛā}, \textquotedblleft{}beautiful\textquotedblright{}, under it (CDIAL 10803) but not \textbf{\mw{रूपाळो}} (\emph{rūpāḷo}) itself.
\end{cousinweb}

\begin{grammarlens}[title={What you've built}]
One more word for this chapter.
\end{grammarlens}

\subsection*{Your turn}
\begin{itemize}
\raggedright
  \item \emph{Say it:} \emph{rūpāḷo}
  \item \emph{Say it:} \emph{rūpāḷo}, once more
  \item \emph{Say it:} say \emph{rūpāḷo}
  \item \emph{Recall:} say the Marwadi for spicy, then the Marwadi for beautiful, then say \emph{rūpāḷo} again
\end{itemize}

\begin{tcolorbox}[breakable,colback=teal!4,colframe=teal!35!black,title={Before you move on}]
What does \mw{रूपाळो} mean? (\textquotedblleft{}Good-looking\textquotedblright{}.) Say it once more.
\end{tcolorbox}
\section[\texorpdfstring{\mw{जवान} (javān)}{जवान (javān)}]{\mw{जवान} (javān) — young}

\label{lesson:MW-C144-javan}



\begin{quote}
Before the new one: say the Marwadi for beautiful, then the Marwadi for good-looking.
\end{quote}

\subsection*{You'll want to know: \mw{जवान}}
\textbf{\mw{जवान}} — \emph{javān} — \textquotedblleft{}young\textquotedblright{}.

\begin{cousinweb}
\textbf{\mw{जवान}} (\emph{javān}) is a loan, usually traced to Persian \emph{javān}, \textquotedblleft{}young\textquotedblright{}.
\end{cousinweb}

\begin{grammarlens}[title={What you've built}]
One more word for this chapter.
\end{grammarlens}

\subsection*{Your turn}
\begin{itemize}
\raggedright
  \item \emph{Say it:} \emph{javān}
  \item \emph{Say it:} \emph{javān}, once more
  \item \emph{Say it:} say \emph{javān}
  \item \emph{Recall:} say the Marwadi for beautiful, then the Marwadi for good-looking, then say \emph{javān} again
\end{itemize}

\begin{tcolorbox}[breakable,colback=teal!4,colframe=teal!35!black,title={Before you move on}]
What does \mw{जवान} mean? (\textquotedblleft{}Young\textquotedblright{}.) Say it once more.
\end{tcolorbox}
\section[\texorpdfstring{\mw{दुबळो} (dubaḷo)}{दुबळो (dubaḷo)}]{\mw{दुबळो} (dubaḷo) — thin, weak}

\label{lesson:MW-C144-dubalo}



\begin{quote}
Before the new one: say the Marwadi for good-looking, then the Marwadi for young.
\end{quote}

\subsection*{You'll want to know: \mw{दुबळो}}
\textbf{\mw{दुबळो}} — \emph{dubaḷo} — \textquotedblleft{}thin, weak\textquotedblright{}.

\begin{cousinweb}
\textbf{\mw{दुबळो}} (\emph{dubaḷo}) is a tadbhava, a word inherited from Sanskrit and reshaped by sound change. It comes from \emph{durbala}, \textquotedblleft{}weak\textquotedblright{}. CDIAL lists Old Marwari \emph{dūbalaü} under it (CDIAL 6438). It is built on \emph{dus-}, \textquotedblleft{}ill\textquotedblright{}, and \emph{bala}, \textquotedblleft{}strength\textquotedblright{}.
\end{cousinweb}

\begin{grammarlens}[title={What you've built}]
One more word for this chapter.
\end{grammarlens}

\subsection*{Your turn}
\begin{itemize}
\raggedright
  \item \emph{Say it:} \emph{dubaḷo}
  \item \emph{Say it:} \emph{dubaḷo}, once more
  \item \emph{Say it:} say \emph{dubaḷo}
  \item \emph{Recall:} say the Marwadi for good-looking, then the Marwadi for young, then say \emph{dubaḷo} again
\end{itemize}

\begin{tcolorbox}[breakable,colback=teal!4,colframe=teal!35!black,title={Before you move on}]
What does \mw{दुबळो} mean? (\textquotedblleft{}Thin, weak\textquotedblright{}.) Say it once more.
\end{tcolorbox}
\section[\texorpdfstring{\mw{ठिगणो} (ṭhigṇo)}{ठिगणो (ṭhigṇo)}]{\mw{ठिगणो} (ṭhigṇo) — short (of a person)}

\label{lesson:MW-C144-thigno}



\begin{quote}
Before the new one: say the Marwadi for young, then the Marwadi for thin.
\end{quote}

\subsection*{You'll want to know: \mw{ठिगणो}}
\textbf{\mw{ठिगणो}} — \emph{ṭhigṇo} — \textquotedblleft{}short (of a person)\textquotedblright{}.

\begin{cousinweb}
\textbf{\mw{ठिगणो}} (\emph{ṭhigṇo}) is a tadbhava, a word inherited and reshaped by sound change. It comes from \emph{*ṭhiṅga}, \textquotedblleft{}defective\textquotedblright{}, a form CDIAL reconstructs from its descendants. CDIAL gives no Marwari form in this entry (CDIAL 5501); the match is through Hindi \emph{ṭhĩgnā} and Gujarati \emph{ṭhĩgṇũ}. CDIAL sets it in a large group of \textquotedblleft{}defective\textquotedblright{} words; Hindi and Gujarati forms mean \textquotedblleft{}dwarfish\textquotedblright{}.
\end{cousinweb}

\begin{grammarlens}[title={What you've built}]
One more word for this chapter.
\end{grammarlens}

\subsection*{Your turn}
\begin{itemize}
\raggedright
  \item \emph{Say it:} \emph{ṭhigṇo}
  \item \emph{Say it:} \emph{ṭhigṇo}, once more
  \item \emph{Say it:} say \emph{ṭhigṇo}
  \item \emph{Recall:} say the Marwadi for young, then the Marwadi for thin, then say \emph{ṭhigṇo} again
\end{itemize}

\begin{tcolorbox}[breakable,colback=teal!4,colframe=teal!35!black,title={Before you move on}]
What does \mw{ठिगणो} mean? (\textquotedblleft{}Short (of a person)\textquotedblright{}.) Say it once more.
\end{tcolorbox}
